\documentclass[11pt]{article}
\usepackage[margin=1in]{geometry}
\usepackage{amsmath,amssymb,amsthm,bm,mathtools}
\usepackage{hyperref}
\usepackage{enumitem}
\usepackage{array}
\usepackage{booktabs}
\usepackage{longtable}
\usepackage{xcolor}

\hypersetup{
  colorlinks=true,
  linkcolor=blue,
  urlcolor=blue,
  citecolor=blue
}

\title{Implementation Summary for Extending PPL-SI with Cross-Validation Selection Events}
\author{}
\date{}

\begin{document}
\maketitle

\section{Purpose of This Note}
This document is a technical implementation note for extending the current \texttt{PPL-SI} codebase with the new \emph{cross-validation-aware} selective inference component described in Section~5 and Appendix~A.8 of the paper \emph{Post-Pretrained Lasso Statistical Inference}. The goal is to help the implementation team translate the current mathematical derivation into code as accurately as possible, while keeping the design aligned with the existing repository.

The main point is:
\begin{itemize}[leftmargin=2em]
    \item the current repository already implements the main PPL-SI pipeline for fixed hyperparameters,
    \item the next new component is to incorporate the \emph{data-driven hyperparameter selection event} induced by cross-validation,
    \item mathematically, this adds two new conditioning sets, $\mathcal{Z}_2$ and $\mathcal{Z}_3$, on top of the existing truncation region logic.
\end{itemize}

\section{What Already Exists}
The paper and repository already provide the base selective-inference framework for:
\begin{enumerate}[leftmargin=2em]
    \item standard Pretrained Lasso,
    \item parameter-only Pretrained Lasso,
    \item D-TransFusion,
    \item divide-and-conquer identification of the truncation region for fixed hyperparameters.
\end{enumerate}

The currently available repository structure is:
\begin{verbatim}
PPL-SI/
|-- ppl_si/
|   |-- __init__.py
|   |-- PPL_SI.py
|   |-- algorithms.py
|   |-- sub_prob.py
|   |-- gen_data.py
|   `-- utils.py
|-- examples/
|   |-- ex1_p_value_PPL.ipynb
|   |-- ex2_p_value_PPL_parameter_only.ipynb
|   |-- ex3_p_value_DTF.ipynb
|   `-- ex4_pivot.ipynb
|-- figures/
`-- README.md
\end{verbatim}

Therefore, the CV extension should be implemented as an \emph{incremental extension} of the existing architecture, not as a separate independent pipeline.

\section{Main Development Target}
For fixed hyperparameters, the current selective inference conditions on
\[
\mathcal{E} = \{\mathcal{A}(\bm{Y}) = \mathcal{A}_{\rm obs},\; \mathcal{Q}(\bm{Y}) = \mathcal{Q}_{\rm obs}\}.
\]
After conditioning on the nuisance statistic, the data are restricted to the line
\[
\bm{Y}(z) = \bm{a} + \bm{b}z,
\]
and the inference problem reduces to characterizing
\[
\mathcal{Z} = \{ z \in \mathbb{R} \mid \mathcal{A}(\bm{Y}(z)) = \mathcal{A}_{\rm obs} \}.
\]

For the new CV-aware version, the conditioning event becomes
\[
\mathcal{E}_{\rm CV} = \{ \mathcal{A}(\bm{Y}) = \mathcal{A}_{\rm obs},\; \mathcal{V}(\bm{Y}) = \mathcal{V}_{\rm obs},\; \tilde{\mathcal{V}}(\bm{Y}) = \tilde{\mathcal{V}}_{\rm obs},\; \mathcal{Q}(\bm{Y}) = \mathcal{Q}_{\rm obs} \}.
\]
Here:
\begin{align*}
    \mathcal{V}:\; \bm{Y} &\mapsto \lambda_{\rm sh} \in \Lambda, \\
    \tilde{\mathcal{V}}:\; \bm{Y} &\mapsto (\rho, \lambda_K) \in \tilde{\Lambda},
\end{align*}
where:
\begin{itemize}[leftmargin=2em]
    \item $\mathcal{V}$ is the Step~1 CV selector for the shared pretraining parameter $\lambda_{\rm sh}$,
    \item $\tilde{\mathcal{V}}$ is the Step~2 CV selector for the target-specific pair $(\rho, \lambda_K)$,
    \item $\mathcal{V}_{\rm obs} = \lambda_{\rm sh}^{\rm obs}$,
    \item $\tilde{\mathcal{V}}_{\rm obs} = (\rho^{\rm obs}, \lambda_K^{\rm obs})$.
\end{itemize}

After conditioning on $\mathcal{Q}(\bm{Y}) = \mathcal{Q}_{\rm obs}$, the conditional data space is
\[
\mathcal{Y}_{\rm CV} = \{ \bm{Y}(z) = \bm{a} + \bm{b}z \mid z \in \mathcal{Z}_{\rm CV} \},
\]
with
\begin{align*}
    \mathcal{Z}_{\rm CV}
    &= \{ z \in \mathbb{R} \mid \mathcal{A}(\bm{Y}) = \mathcal{A}_{\rm obs},\; \mathcal{V}(\bm{Y}) = \mathcal{V}_{\rm obs},\; \tilde{\mathcal{V}}(\bm{Y}) = \tilde{\mathcal{V}}_{\rm obs} \} \\
    &= \underbrace{\{ z \in \mathbb{R} \mid \mathcal{A}(\bm{Y}) = \mathcal{A}_{\rm obs} \}}_{\mathcal{Z}_1}
    \cap
    \underbrace{\{ z \in \mathbb{R} \mid \mathcal{V}(\bm{Y}) = \mathcal{V}_{\rm obs} \}}_{\mathcal{Z}_2}
    \cap
    \underbrace{\{ z \in \mathbb{R} \mid \tilde{\mathcal{V}}(\bm{Y}) = \tilde{\mathcal{V}}_{\rm obs} \}}_{\mathcal{Z}_3}.
\end{align*}

The existing code already targets $\mathcal{Z}_1$. The new development is to construct $\mathcal{Z}_2$ and $\mathcal{Z}_3$, then intersect them with $\mathcal{Z}_1$.

\section{How the Existing Fixed-Hyperparameter Pipeline Should Be Reused}
The current implementation for fixed hyperparameters should remain the base layer.

\subsection*{Existing role of $\mathcal{Z}_1$}
$\mathcal{Z}_1$ is the feature-selection event of the final Pretrained Lasso output. The repository already implements the divide-and-conquer logic for this part through the decomposition into:
\[
\mathcal{Z}_{u,v,t} = \mathcal{Z}_u \cap \mathcal{Z}_v \cap \mathcal{Z}_t,
\]
where:
\begin{itemize}[leftmargin=2em]
    \item $\mathcal{Z}_u$ is the Step~1 pretraining active-set/sign event,
    \item $\mathcal{Z}_v$ is the Step~2 fine-tuning active-set/sign event,
    \item $\mathcal{Z}_t$ is the final aggregated active-set/sign event.
\end{itemize}

These are all characterized by linear inequalities in $z$ and should stay unchanged.

\subsection*{New role of $\mathcal{Z}_2$ and $\mathcal{Z}_3$}
The CV extension does \emph{not} replace the existing line-search / interval-search engine. Instead, it adds extra filters:
\begin{itemize}[leftmargin=2em]
    \item $\mathcal{Z}_2$ ensures that along the line $\bm{Y}(z)$, Step~1 CV still chooses the observed winner $\lambda_{\rm sh}^{\rm obs}$,
    \item $\mathcal{Z}_3$ ensures that along the line $\bm{Y}(z)$, Step~2 CV still chooses the observed winner $(\rho^{\rm obs}, \lambda_K^{\rm obs})$.
\end{itemize}

So the final truncation region becomes
\[
\mathcal{Z}_{\rm CV} = \mathcal{Z}_1 \cap \mathcal{Z}_2 \cap \mathcal{Z}_3.
\]

\section{Implementation Assumptions}
The implementation should explicitly adopt the following assumptions:
\begin{enumerate}[leftmargin=2em]
    \item The candidate grids $\Lambda$ and $\tilde{\Lambda}$ are fixed beforehand.
    \item The fold split used by cross-validation is fixed beforehand.
    \item Tie-breaking for CV minimizers is deterministic.
    \item The Step~1 CV event and Step~2 CV event follow the same sequential two-stage logic as in the Pretrained Lasso implementation, rather than an expensive fully nested global CV of the entire pipeline in one step.
\end{enumerate}

These assumptions should be documented directly in code comments and in the example notebook for the CV extension.

\section{Step~1 CV Event: Characterization of $\mathcal{Z}_2$}

\subsection{Training/validation split notation}
For one validation split, rewrite the pooled data as
\[
\{X, \bm{Y}\} = \left\{ \begin{pmatrix} X_{\rm train} \\ X_{\rm val} \end{pmatrix},\; \begin{pmatrix} \bm{Y}_{\rm train} \\ \bm{Y}_{\rm val} \end{pmatrix} \right\}.
\]
Under the affine line parameterization,
\[
\bm{Y}_{\rm train}(z) = \bm{a}_{\rm train} + \bm{b}_{\rm train} z,
\qquad
\bm{Y}_{\rm val}(z) = \bm{a}_{\rm val} + \bm{b}_{\rm val} z.
\]

For any $\lambda_{\rm sh} \in \Lambda$, the Step~1 estimator on the training split is
\[
\hat{\bm\beta}^{\rm sh}_{\lambda_{\rm sh}}(z)
=
\operatorname*{argmin}_{\bm\beta \in \mathbb{R}^{p^*}}
\frac{1}{2}\left\| \bm{Y}_{\rm train}(z) - X_{\rm train}\bm\beta \right\|_2^2
+
\lambda_{\rm sh}\sum_{j=2}^{p^*} |\beta_j|.
\]

The validation error is
\[
E_{\lambda_{\rm sh}}(z)
=
\frac{1}{2}
\left\| \bm{Y}_{\rm val}(z) - X_{\rm val}\hat{\bm\beta}^{\rm sh}_{\lambda_{\rm sh}}(z) \right\|_2^2.
\]

The Step~1 CV event is
\[
\mathcal{Z}_2
=
\left\{ z \in \mathbb{R} \, \middle| \,
\bar{E}_{\lambda_{\rm sh}^{\rm obs}}(z) \le \bar{E}_{\lambda_{\rm sh}}(z)
\;\text{for all } \lambda_{\rm sh} \in \Lambda
\right\}.
\]

\subsection{Piecewise-linear Step~1 estimator}
For each candidate $\lambda_{\rm sh}$, the estimator $\hat{\bm\beta}^{\rm sh}_{\lambda_{\rm sh}}(z)$ is piecewise linear in $z$. Concretely, on each interval $[z_u, z_{u+1})$ where the active set and sign pattern remain fixed,
\[
\hat{\bm\beta}^{\rm sh}_{\lambda_{\rm sh}}(z)
=
\bm{g}_{\lambda_{\rm sh}}^{u} + \bm{h}_{\lambda_{\rm sh}}^{u} z.
\]
The coefficients are
\begin{align*}
\bm{g}_{\lambda_{\rm sh}}^{u}
&=
E_{\mathcal{C}_u}
\left(X_{\rm train,\mathcal{C}_u}^{\top}X_{\rm train,\mathcal{C}_u}\right)^{-1}
\left(X_{\rm train,\mathcal{C}_u}^{\top}\bm{a}_{\rm train} - \tilde{\bm w}_{\mathcal{C}_u} \circ \mathcal{S}_{\mathcal{C}_u}\right),
\\
\bm{h}_{\lambda_{\rm sh}}^{u}
&=
E_{\mathcal{C}_u}
\left(X_{\rm train,\mathcal{C}_u}^{\top}X_{\rm train,\mathcal{C}_u}\right)^{-1}
X_{\rm train,\mathcal{C}_u}^{\top}\bm{b}_{\rm train}.
\end{align*}

Operationally, this means the implementation should already be able to recover interval-wise affine coefficients from the existing sub-problem machinery or from a dedicated path routine for the Step~1 training-fold Lasso.

\subsection{Piecewise-quadratic validation error}
On each interval $[z_u, z_{u+1})$,
\begin{align*}
E_{\lambda_{\rm sh}}(z)
&=
\frac{1}{2}
\left\|
(\bm{a}_{\rm val} - X_{\rm val}\bm{g}_{\lambda_{\rm sh}}^{u})
+
(\bm{b}_{\rm val} - X_{\rm val}\bm{h}_{\lambda_{\rm sh}}^{u}) z
\right\|_2^2
\\
&=
A_{\lambda_{\rm sh}}^{u} z^2 + B_{\lambda_{\rm sh}}^{u} z + C_{\lambda_{\rm sh}}^{u},
\end{align*}
where
\begin{align*}
A_{\lambda_{\rm sh}}^{u}
&=
\frac{1}{2}
(\bm{b}_{\rm val} - X_{\rm val}\bm{h}_{\lambda_{\rm sh}}^{u})^{\top}
(\bm{b}_{\rm val} - X_{\rm val}\bm{h}_{\lambda_{\rm sh}}^{u}),
\\
B_{\lambda_{\rm sh}}^{u}
&=
(\bm{a}_{\rm val} - X_{\rm val}\bm{g}_{\lambda_{\rm sh}}^{u})^{\top}
(\bm{b}_{\rm val} - X_{\rm val}\bm{h}_{\lambda_{\rm sh}}^{u}),
\\
C_{\lambda_{\rm sh}}^{u}
&=
\frac{1}{2}
(\bm{a}_{\rm val} - X_{\rm val}\bm{g}_{\lambda_{\rm sh}}^{u})^{\top}
(\bm{a}_{\rm val} - X_{\rm val}\bm{g}_{\lambda_{\rm sh}}^{u}).
\end{align*}

\subsection{Extension to $M$-fold CV}
For $M$ folds, define the mean validation error
\[
\bar{E}_{\lambda_{\rm sh}}(z) = \frac{1}{M} \sum_{m=1}^{M} E_{\lambda_{\rm sh}}^{(m)}(z).
\]
The implementation must do the following for each candidate $\lambda_{\rm sh}$:
\begin{enumerate}[leftmargin=2em]
    \item For each fold $m$, compute the breakpoints of the piecewise-linear estimator and therefore the breakpoints of $E_{\lambda_{\rm sh}}^{(m)}(z)$.
    \item Merge all fold-specific breakpoints into a common interval partition.
    \item On each common interval $I_r$, represent each fold-specific validation error by a quadratic
    \[
    E_{\lambda_{\rm sh}}^{(m)}(z) = A_{\lambda_{\rm sh}}^{(m,r)} z^2 + B_{\lambda_{\rm sh}}^{(m,r)} z + C_{\lambda_{\rm sh}}^{(m,r)}.
    \]
    \item Average the coefficients:
    \begin{align*}
        \bar{A}_{\lambda_{\rm sh}}^{(r)} &= \frac{1}{M}\sum_{m=1}^M A_{\lambda_{\rm sh}}^{(m,r)}, \\
        \bar{B}_{\lambda_{\rm sh}}^{(r)} &= \frac{1}{M}\sum_{m=1}^M B_{\lambda_{\rm sh}}^{(m,r)}, \\
        \bar{C}_{\lambda_{\rm sh}}^{(r)} &= \frac{1}{M}\sum_{m=1}^M C_{\lambda_{\rm sh}}^{(m,r)}.
    \end{align*}
\end{enumerate}

Then, for every competitor $\lambda_{\rm sh} \in \Lambda$, construct the pairwise winner inequality against the observed winner $\lambda_{\rm sh}^{\rm obs}$:
\[
\bar{E}_{\lambda_{\rm sh}^{\rm obs}}(z) - \bar{E}_{\lambda_{\rm sh}}(z) \le 0.
\]
On each common interval this becomes
\[
A_{\lambda_{\rm sh}} z^2 + B_{\lambda_{\rm sh}} z + C_{\lambda_{\rm sh}} \le 0,
\]
with
\begin{align*}
A_{\lambda_{\rm sh}} &= \bar{A}_{\lambda_{\rm sh}^{\rm obs}} - \bar{A}_{\lambda_{\rm sh}}, \\
B_{\lambda_{\rm sh}} &= \bar{B}_{\lambda_{\rm sh}^{\rm obs}} - \bar{B}_{\lambda_{\rm sh}}, \\
C_{\lambda_{\rm sh}} &= \bar{C}_{\lambda_{\rm sh}^{\rm obs}} - \bar{C}_{\lambda_{\rm sh}}.
\end{align*}

Therefore, operationally,
\[
\mathcal{Z}_2 = \bigcap_{\lambda_{\rm sh} \in \Lambda}
\left\{ z \in \mathbb{R} \mid A_{\lambda_{\rm sh}} z^2 + B_{\lambda_{\rm sh}} z + C_{\lambda_{\rm sh}} \le 0 \right\},
\]
with the understanding that these coefficients are evaluated interval by interval on the merged partition.

\section{Step~2 CV Event: Characterization of $\mathcal{Z}_3$}

\subsection{Target-domain split notation}
The Step~2 CV event acts only on the target domain $K$. Rewrite the target-domain data as
\[
\{X^{(K)}, \bm{Y}^{(K)}\}
=
\left\{
\begin{pmatrix} X^{(K)}_{\rm train} \\ X^{(K)}_{\rm val} \end{pmatrix},
\begin{pmatrix} \bm{Y}^{(K)}_{\rm train} \\ \bm{Y}^{(K)}_{\rm val} \end{pmatrix}
\right\}.
\]
The parameterized responses are
\[
\bm{Y}^{(K)}_{\rm train}(z) = \bm{a}^{(K)}_{\rm train} + \bm{b}^{(K)}_{\rm train} z,
\qquad
\bm{Y}^{(K)}_{\rm val}(z) = \bm{a}^{(K)}_{\rm val} + \bm{b}^{(K)}_{\rm val} z.
\]

Define $\Phi = (\rho, \lambda_K) \in \tilde{\Lambda}$. The Step~2 individual estimator is
\[
\hat{\bm\beta}_{\Phi}^{(K){\rm indiv}}(z)
=
\operatorname*{argmin}_{\bm\beta \in \mathbb{R}^{p^*}}
\frac{1}{2}
\left\|
\bm{Y}_{\rm train}^{(K)}(z)
-
(1-\rho)X_{\rm train}^{(K)}\hat{\bm\beta}^{\rm sh}(z)
-
X_{\rm train}^{(K)}\bm\beta
\right\|_2^2
+
\lambda_K \sum_{j=2}^{p^*} a_j |\beta_j|.
\]
The final estimator is
\[
\hat{\bm\beta}_{\Phi}^{(K)}(z)
=
(1-\rho)\hat{\bm\beta}^{\rm sh}(z) + \hat{\bm\beta}_{\Phi}^{(K){\rm indiv}}(z).
\]

The validation error is
\[
\tilde{E}_{\Phi}(z)
=
\frac{1}{2}
\left\| \bm{Y}^{(K)}_{\rm val}(z) - X^{(K)}_{\rm val}\hat{\bm\beta}_{\Phi}^{(K)}(z) \right\|_2^2.
\]

The Step~2 CV event is
\[
\mathcal{Z}_3
=
\left\{ z \in \mathbb{R} \, \middle| \,
\bar{\tilde{E}}_{\Phi^{\rm obs}}(z) \le \bar{\tilde{E}}_{\Phi}(z)
\;\text{for all } \Phi \in \tilde{\Lambda}
\right\}.
\]

\subsection{Important implementation rule}
When constructing $\mathcal{Z}_3$, the Step~1 shared parameter is fixed at the observed CV-selected winner:
\[
\lambda_{\rm sh} = \lambda_{\rm sh}^{\rm obs}.
\]
This is critical. The Step~2 CV event is defined conditional on the observed Step~1 winner, exactly matching the sequential two-stage CV logic.

\subsection{Piecewise-linear final estimator}
The Step~1 pretrained estimator is piecewise linear in $z$, and the Step~2 individual estimator is also piecewise linear in $z$. Therefore, on each interval formed by merging both breakpoint sets,
\[
\hat{\bm\beta}_{\Phi}^{(K)}(z) = \tilde{\bm g}_{\Phi}^{o} + \tilde{\bm h}_{\Phi}^{o} z.
\]
Operationally, each interval index $o$ corresponds to the simultaneous stability of:
\begin{itemize}[leftmargin=2em]
    \item the Step~1 active set/sign pattern,
    \item the Step~2 individual active set/sign pattern.
\end{itemize}

\subsection{Piecewise-quadratic Step~2 validation error}
On each interval $[z_o, z_{o+1})$,
\begin{align*}
\tilde{E}_{\Phi}(z)
&=
\frac{1}{2}
\left\|
(\bm{a}^{(K)}_{\rm val} - X^{(K)}_{\rm val}\tilde{\bm g}_{\Phi}^{o})
+
(\bm{b}^{(K)}_{\rm val} - X^{(K)}_{\rm val}\tilde{\bm h}_{\Phi}^{o})z
\right\|_2^2
\\
&=
\tilde{A}_{\Phi}^{o} z^2 + \tilde{B}_{\Phi}^{o} z + \tilde{C}_{\Phi}^{o},
\end{align*}
where
\begin{align*}
\tilde{A}_{\Phi}^{o}
&=
\frac{1}{2}
(\bm{b}^{(K)}_{\rm val} - X^{(K)}_{\rm val}\tilde{\bm h}_{\Phi}^{o})^{\top}
(\bm{b}^{(K)}_{\rm val} - X^{(K)}_{\rm val}\tilde{\bm h}_{\Phi}^{o}),
\\
\tilde{B}_{\Phi}^{o}
&=
(\bm{a}^{(K)}_{\rm val} - X^{(K)}_{\rm val}\tilde{\bm g}_{\Phi}^{o})^{\top}
(\bm{b}^{(K)}_{\rm val} - X^{(K)}_{\rm val}\tilde{\bm h}_{\Phi}^{o}),
\\
\tilde{C}_{\Phi}^{o}
&=
\frac{1}{2}
(\bm{a}^{(K)}_{\rm val} - X^{(K)}_{\rm val}\tilde{\bm g}_{\Phi}^{o})^{\top}
(\bm{a}^{(K)}_{\rm val} - X^{(K)}_{\rm val}\tilde{\bm g}_{\Phi}^{o}).
\end{align*}

For $M$-fold CV, define
\[
\bar{\tilde{E}}_{\Phi}(z) = \frac{1}{M}\sum_{m=1}^{M} \tilde{E}_{\Phi}^{(m)}(z).
\]
As in Step~1, implementation must:
\begin{enumerate}[leftmargin=2em]
    \item compute fold-specific breakpoint sets,
    \item merge them into a common interval partition,
    \item express each fold error as a quadratic on each common interval,
    \item average coefficients across folds.
\end{enumerate}

Then, for each competitor $\Phi \in \tilde{\Lambda}$,
\[
\bar{\tilde{E}}_{\Phi^{\rm obs}}(z) - \bar{\tilde{E}}_{\Phi}(z) \le 0.
\]
On each common interval this becomes
\[
\tilde{A}_{\Phi} z^2 + \tilde{B}_{\Phi} z + \tilde{C}_{\Phi} \le 0,
\]
with
\begin{align*}
\tilde{A}_{\Phi} &= \bar{\tilde{A}}_{\Phi^{\rm obs}} - \bar{\tilde{A}}_{\Phi}, \\
\tilde{B}_{\Phi} &= \bar{\tilde{B}}_{\Phi^{\rm obs}} - \bar{\tilde{B}}_{\Phi}, \\
\tilde{C}_{\Phi} &= \bar{\tilde{C}}_{\Phi^{\rm obs}} - \bar{\tilde{C}}_{\Phi}.
\end{align*}

Therefore,
\[
\mathcal{Z}_3 = \bigcap_{\Phi \in \tilde{\Lambda}}
\left\{ z \in \mathbb{R} \mid \tilde{A}_{\Phi} z^2 + \tilde{B}_{\Phi} z + \tilde{C}_{\Phi} \le 0 \right\},
\]
again interpreted interval by interval after merging breakpoints.

\section{Selective p-Value with CV}
Once $\mathcal{Z}_{\rm CV}$ is obtained, inference proceeds exactly as in the fixed-hyperparameter case, except the truncation region is now tighter. Let
\[
Z = \eta_j^{\top}\bm{Y},
\qquad
Z_{\rm obs} = \eta_j^{\top}\bm{Y}_{\rm obs}.
\]
Then the CV-aware selective p-value is computed from the truncated normal law restricted to $\mathcal{Z}_{\rm CV}$:
\[
p^{\rm selective, CV}_j
=
\mathbb{P}_{H_{0,j}}
\left(
|Z| \ge |Z_{\rm obs}|
\;\middle|\;
Z \in \mathcal{Z}_{\rm CV}
\right).
\]

So from a code perspective, the only new work is to replace the old truncation region input $\mathcal{Z}$ with the new one:
\[
\mathcal{Z}_{\rm CV} = \mathcal{Z}_1 \cap \mathcal{Z}_2 \cap \mathcal{Z}_3.
\]
The final p-value engine itself should not need fundamental redesign.

\section{Recommended Code-Level Design}
A practical way to implement the CV extension inside the current repository is to add a small set of new reusable functions.

\subsection{Recommended new components}
\begin{enumerate}[leftmargin=2em]
    \item \texttt{cv\_utils.py} or new functions inside \texttt{utils.py}
    \begin{itemize}[leftmargin=2em]
        \item fold generation,
        \item breakpoint merging,
        \item interval alignment,
        \item quadratic coefficient averaging,
        \item interval intersection utilities for quadratic constraints.
    \end{itemize}

    \item New Step~1 CV path routine
    \begin{itemize}[leftmargin=2em]
        \item input: $(X, \bm a, \bm b, \Lambda, \text{folds})$,
        \item output: interval-wise quadratic representation of $\bar{E}_{\lambda_{\rm sh}}(z)$ for all $\lambda_{\rm sh} \in \Lambda$.
    \end{itemize}

    \item New Step~2 CV path routine
    \begin{itemize}[leftmargin=2em]
        \item input: $(X^{(K)}, \bm a^{(K)}, \bm b^{(K)}, \lambda_{\rm sh}^{\rm obs}, \tilde{\Lambda}, \text{folds})$,
        \item output: interval-wise quadratic representation of $\bar{\tilde{E}}_{\Phi}(z)$ for all $\Phi \in \tilde{\Lambda}$.
    \end{itemize}

    \item New truncation-region constructor
    \begin{itemize}[leftmargin=2em]
        \item build $\mathcal{Z}_2$ from Step~1 quadratic inequalities,
        \item build $\mathcal{Z}_3$ from Step~2 quadratic inequalities,
        \item intersect with the already existing $\mathcal{Z}_1$.
    \end{itemize}

    \item New example notebook
    \begin{itemize}[leftmargin=2em]
        \item suggested name: \texttt{ex5\_p\_value\_PPL\_CV.ipynb}.
    \end{itemize}
\end{enumerate}

\subsection{Where each part likely belongs}
\begin{center}
\begin{tabular}{p{0.28\textwidth} p{0.62\textwidth}}
\toprule
\textbf{File} & \textbf{Recommended responsibility} \\
\midrule
\texttt{algorithms.py} & estimator calls for Step~1 and Step~2 across CV folds and hyperparameter grids \\
\texttt{sub\_prob.py} & existing fixed-hyperparameter interval logic; may also host interval-wise stability helpers reused by CV \\
\texttt{utils.py} & interval merging, quadratic coefficient bookkeeping, winner-inequality construction \\
\texttt{PPL\_SI.py} & top-level orchestration: compute $\mathcal{Z}_1$, $\mathcal{Z}_2$, $\mathcal{Z}_3$, intersect them, then evaluate the selective p-value \\
\bottomrule
\end{tabular}
\end{center}

\section{Concrete End-to-End Implementation Flow}
The final implementation should follow this order:
\begin{enumerate}[leftmargin=2em]
    \item Run the observed-data pipeline once to obtain:
    \begin{itemize}[leftmargin=2em]
        \item $\mathcal{A}_{\rm obs}$,
        \item $\lambda_{\rm sh}^{\rm obs}$,
        \item $(\rho^{\rm obs}, \lambda_K^{\rm obs})$.
    \end{itemize}

    \item For each selected feature $j \in \mathcal{A}_{\rm obs}$:
    \begin{enumerate}[leftmargin=2em]
        \item compute $\eta_j$, $\bm a$, $\bm b$,
        \item build $\mathcal{Z}_1$ using the existing divide-and-conquer algorithm,
        \item build $\mathcal{Z}_2$ using Step~1 CV quadratic inequalities,
        \item build $\mathcal{Z}_3$ using Step~2 CV quadratic inequalities,
        \item intersect all three sets,
        \item evaluate the truncated normal tail probability on the resulting region.
    \end{enumerate}
\end{enumerate}

\section{Validation Checklist for the Implementation Team}
The implementation should be considered correct only if all items below pass.

\subsection*{A. Structural checks}
\begin{itemize}[leftmargin=2em]
    \item If the hyperparameter grids contain only one value each, the CV-aware code should reduce to the existing fixed-hyperparameter code.
    \item $\mathcal{Z}_{\rm CV}$ must always be a subset of $\mathcal{Z}_1$.
    \item The observed $z_{\rm obs}$ must belong to the final truncation region.
\end{itemize}

\subsection*{B. Numerical checks}
\begin{itemize}[leftmargin=2em]
    \item For each fold and candidate hyperparameter, the validation error computed directly should match the interval-wise quadratic evaluation.
    \item On random test points inside an interval, the value of the fitted quadratic should match the brute-force validation error up to numerical tolerance.
    \item Pairwise winner inequalities should be satisfied at $z_{\rm obs}$ for the observed winner.
\end{itemize}

\subsection*{C. Statistical checks}
\begin{itemize}[leftmargin=2em]
    \item Under the null, the pivot should be approximately uniform.
    \item FPR should remain controlled at the nominal level.
    \item Compared with naive inference, the CV-aware method should avoid false p-value inflation.
\end{itemize}

\section{Recommended Minimal Deliverables}
The CV extension should ideally be delivered with the following items:
\begin{enumerate}[leftmargin=2em]
    \item code implementation in the \texttt{ppl\_si} package,
    \item one example notebook for CV-aware p-values,
    \item one notebook or script for pivot validation,
    \item one notebook or script for FPR/TPR comparison,
    \item short README update describing how to enable CV-aware selective inference.
\end{enumerate}

\section{Final Summary}
The repository already contains the main PPL-SI machinery for fixed hyperparameters. The new task is specifically to implement the CV-aware extension by adding the two extra truncation-region components:
\[
\mathcal{Z}_{\rm CV} = \mathcal{Z}_1 \cap \mathcal{Z}_2 \cap \mathcal{Z}_3.
\]

The practical implementation should follow the same high-level principle in both Step~1 and Step~2:
\begin{enumerate}[leftmargin=2em]
    \item along the conditioned line $\bm{Y}(z) = \bm a + \bm b z$, represent estimators as piecewise-linear functions of $z$,
    \item turn validation errors into piecewise-quadratic functions of $z$,
    \item compare each candidate hyperparameter against the observed winner,
    \item convert these winner conditions into quadratic inequalities,
    \item intersect those CV constraints with the existing selection-event region,
    \item compute the selective p-value on the final truncation region.
\end{enumerate}

This is the key new development direction for the codebase.

\vspace{1em}
\noindent
Repository reference: \href{https://github.com/DAIR-Group/PPL-SI}{DAIR-Group/PPL-SI}.

\end{document}
